% ============================================================
%  IEEE JBHI — Full Paper
%  Title: PrenatalSentinel: A Privacy-Preserving Federated
%         Learning Framework for Autism Risk Prediction from
%         Multimodal Prenatal Purchase Signals
%  Author: Kapil Kumar Reddy Poreddy (poreddykapil@ieee.org)
%  Repository: github.com/learningdebunked/responsible_ai_pediatrichealth
% ============================================================

\documentclass[journal]{IEEEtran}

\usepackage[T1]{fontenc}
\usepackage[utf8]{inputenc}
\usepackage{amsmath,amssymb,amsthm}
\usepackage{graphicx}
\usepackage{booktabs}
\usepackage{multirow}
\usepackage{array}
\usepackage{xcolor}
\usepackage{hyperref}
\usepackage{cite}
\usepackage{algorithm}
\usepackage{algorithmic}
\usepackage{balance}
\usepackage{url}

\newtheorem{definition}{Definition}
\newtheorem{proposition}{Proposition}

\newcommand{\eg}{\textit{e.g.}}
\newcommand{\ie}{\textit{i.e.}}
\newcommand{\etal}{\textit{et al.}}
\newcommand{\eps}{\varepsilon}
\newcommand{\Lc}{\mathcal{L}}
\newcommand{\Dc}{\mathcal{D}}
\newcommand{\Mc}{\mathcal{M}}

\title{PrenatalSentinel: A Privacy-Preserving Federated Learning
       Framework for Autism Risk Prediction from Multimodal
       Prenatal Purchase Signals}

\author{Kapil~Kumar~Reddy~Poreddy,~\IEEEmembership{Member,~IEEE}
\thanks{Manuscript received June 2025.}
\thanks{K.~K.~R.~Poreddy is an AI researcher and IEEE member.
        E-mail: poreddykapil@ieee.org}
\thanks{Code and dataset: \url{https://github.com/learningdebunked/responsible_ai_pediatrichealth}}
}

\markboth{IEEE Journal of Biomedical and Health Informatics}%
{Poreddy: PrenatalSentinel --- Federated Autism Risk Prediction}

\begin{document}
\maketitle

% ── Abstract ──────────────────────────────────────────────
\begin{abstract}
Autism Spectrum Disorder (ASD) affects 1 in 31 children in the United
States as of 2025 surveillance data, yet the average age of diagnosis
remains 4--5 years, well beyond the optimal early-intervention window
before age 3. We present \textbf{PrenatalSentinel}, a privacy-preserving
federated learning framework that integrates three complementary prenatal
signals derived from publicly available retail and public health
datasets: (1)~grocery food purchase patterns operationalized as a
\textit{food-as-medicine} signal, (2)~over-the-counter (OTC) supplement
and medication purchase patterns, and (3)~USDA Food Desert Index and
socioeconomic status (SES) stratification from the National Survey of
Children's Health (NSCH). Using a synthetic case-enriched cohort of
$N=5{,}000$ records with 13 engineered features, calibrated to known
epidemiological associations, we demonstrate that a federated gradient
aggregation model with $(\eps,\delta)$-differential privacy achieves
AUC~$=~0.685$ at strict privacy budget $\eps=1$, compared to a
centralized logistic regression baseline of AUC~$=~0.672$. Our equity
audit reveals that food-desert zip codes exhibit systematically lower
AUC (0.607 vs.~0.628), quantifying an AI detection gap for low-income
populations. A signal ablation study confirms that adding SES and food
desert features (F+M+E configuration) yields the highest AUC of 0.650,
a 2.4\% improvement over food signals alone. PrenatalSentinel addresses
five open research gaps and proposes a national ASD surveillance
architecture operable through existing grocery and pharmacy loyalty
programs without exposing individual health records.
\end{abstract}

\begin{IEEEkeywords}
Autism spectrum disorder, federated learning, differential privacy,
prenatal nutrition, food as medicine, OTC medications, food desert,
socioeconomic status, early detection, responsible AI, pediatric health
\end{IEEEkeywords}

% ============================================================
\section{Introduction}
\label{sec:intro}

\IEEEPARstart{A}{utism} Spectrum Disorder (ASD) is a neurodevelopmental
condition affecting an estimated 1 in 31 children aged 8 years in the
United States as of the most recent Centers for Disease Control and
Prevention (CDC) Autism and Developmental Disabilities Monitoring
(ADDM) Network surveillance data~\cite{cdc2025addm}, representing a
fivefold increase from 1 in 150 in the year 2000. The cumulative
societal cost of ASD in the United States now exceeds \$461 billion
annually, encompassing healthcare, special education, and lost
productivity~\cite{buescher2014costs}. Despite this scale, the average
age of ASD diagnosis in the United States remains 4--5 years~\cite{cdc2025addm},
substantially beyond the critical neurodevelopmental window before age
3 during which behavioral interventions produce the greatest long-term
gains~\cite{dawson2010randomized}.

The prenatal period represents a biologically plausible upstream window
for ASD risk surveillance. Converging evidence from large prospective
cohorts---including the Norwegian Mother, Father, and Child Cohort
(MoBa) and the Avon Longitudinal Study of Parents and Children
(ALSPAC)---demonstrates that maternal dietary patterns during pregnancy,
particularly adherence to nutrient-dense diets rich in folate,
omega-3 fatty acids, iron, and vitamin D, are associated with
measurably lower offspring ASD risk~\cite{devilbiss2017prenatal,
lyall2021maternal}. Periconceptional folic acid supplementation has
been identified as a protective factor across multiple meta-analyses,
with a pooled relative risk reduction of 23--38\% in high-adherence
groups~\cite{guo2020prenatal}. Conversely, food insecurity and
residence in USDA-designated food deserts are associated with reduced
dietary quality and systematically worse developmental outcomes,
disproportionately affecting low-income and minority
families~\cite{panjwani2021food}.

A critical gap exists between these epidemiological findings and
actionable population-scale surveillance. Existing ASD screening
tools---including the Modified Checklist for Autism in Toddlers
(M-CHAT) and the Quantitative Checklist for Autism in Toddlers
(Q-CHAT-10)---are applied reactively, typically after behavioral
concerns emerge at age 2--3. No system currently leverages the
\textit{prenatal behavioral signal} embedded in grocery and pharmacy
purchase records to generate prospective, population-scale ASD
early-warning scores.

Grocery and pharmacy loyalty programs represent an underutilized
epidemiological resource. The Instacart Market Basket Analysis
dataset~\cite{instacart2017} captures over 3 million orders from
206,209 users across 49,688 products, with product-level aisle
classifications that enable operationalization of dietary quality
metrics. Crucially, these purchase records are generated \textit{before}
any ASD diagnosis, creating a prospective signal window. The challenge
is twofold: (1)~extracting clinically meaningful nutritional features
from raw purchase data, and (2)~doing so at population scale without
violating individual privacy.

Federated learning (FL)~\cite{mcmahan2017communication} addresses the
privacy challenge by training models across distributed retail nodes
without sharing raw transaction data. When combined with
$(\eps,\delta)$-differential privacy~\cite{dwork2014algorithmic},
federated models provide formal, mathematically bounded privacy
guarantees. However, to date, no federated learning system has
integrated grocery purchase signals, OTC supplement purchase patterns,
and food desert index data into a unified prenatal ASD risk framework.

A further dimension of concern is algorithmic equity. If predictive
models are trained predominantly on data from high-SES, food-secure
populations, their performance may degrade precisely for the
low-income families at highest risk---an AI detection gap that
exacerbates existing health disparities. The USDA Food Access Research
Atlas (FARA)~\cite{usdafara2019} provides census-tract-level food
desert designations that enable explicit equity auditing of model
performance across socioeconomic strata.

\textbf{Contributions.} The main contributions of this paper are:
\begin{itemize}
  \item \textbf{C1 --- Novel signal}: First systematic framework for
        using retail grocery and pharmacy OTC purchase data as a
        prospective prenatal ASD risk biomarker at population scale.
  \item \textbf{C2 --- Food-as-medicine operationalization}: A
        purchase-to-nutrient mapping pipeline linking retail product
        categories to USDA FoodData Central nutrient profiles,
        generating six food-quality scores per user.
  \item \textbf{C3 --- Multimodal fusion}: A three-signal federated
        architecture combining food (F), OTC (M), and SES/food desert
        (E) signals into a unified 13-feature ASD risk vector.
  \item \textbf{C4 --- Privacy-preserving}: $(\eps,\delta)$-DP
        guarantees across federated retail nodes with no individual
        data sharing; AUC$=0.685$ at $\eps=1$.
  \item \textbf{C5 --- Equity audit}: Quantification of the ASD
        detection gap across five SES quintiles and food desert
        subgroups using the USDA FARA, revealing a 10.2\% AUC gap
        between high and low food desert populations.
\end{itemize}

The remainder of this paper is organized as follows.
Section~\ref{sec:related} reviews related work across four dimensions.
Section~\ref{sec:datasets} describes the datasets and data linkage
pipeline. Section~\ref{sec:method} presents the feature engineering
and federated learning methodology. Section~\ref{sec:experiments}
details the experimental setup. Section~\ref{sec:results} presents
results. Section~\ref{sec:discussion} discusses implications and
limitations. Section~\ref{sec:ethics} addresses responsible AI and
ethics. Section~\ref{sec:conclusion} concludes.

\begin{figure*}[t]
  \centering
  \includegraphics[width=\textwidth]{prenatal_sentinel_architecture.png}
  \caption{PrenatalSentinel system architecture. Three publicly available
           data streams (grocery purchase records, OTC/supplement purchases,
           and USDA FARA + NSCH) are feature-engineered into a 13-dimensional
           vector per zip code. Five retail client nodes train local gradient
           boosting models and submit clipped gradients to a federated
           aggregation server with Gaussian differential privacy ($\eps,\delta$).
           The aggregated model produces ASD risk scores stratified by SES
           quintile for equity-aware surveillance. No raw data leaves retail
           nodes.}
  \label{fig:arch}
\end{figure*}

% ============================================================
\section{Background and Related Work}
\label{sec:related}

\subsection{Prenatal Nutrition and ASD Risk}

Large prospective birth cohort studies have established associations
between prenatal dietary patterns and offspring ASD risk. DeVilbiss
\etal~\cite{devilbiss2017prenatal} analyzed data from the MoBa and
ALSPAC cohorts using causally informed methods including Mendelian
randomization and found that adherence to a healthy prenatal dietary
pattern (HPDP) was associated with lower probability of offspring ASD
diagnoses and autism-related traits. Lyall \etal~\cite{lyall2021maternal}
conducted a systematic synthesis of maternal nutritional factors and
ASD, identifying folic acid and prenatal vitamins as the most
consistently protective dietary exposures.

Schmidt \etal~\cite{schmidt2012prenatal} reported in the CHARGE study
that prenatal vitamin use during the first month of pregnancy was
associated with reduced ASD risk in children, particularly in mothers
and children with metabolic variants affecting folate processing.
Guo \etal~\cite{guo2020prenatal} conducted a meta-analysis of 22
studies and found that periconceptional folic acid supplementation
was associated with a 23\% reduction in ASD risk (pooled OR 0.77,
95\% CI 0.63--0.93). Friel \etal~\cite{friel2021prenatal} further
examined the role of omega-3 polyunsaturated fatty acids (PUFAs),
finding that lower midpregnancy total PUFA levels were associated with
ASD with intellectual disability in a California case-control study.

Importantly, these cohort studies rely on biomarker measurements or
detailed dietary recall instruments that are not available at population
scale. PrenatalSentinel addresses this gap by operationalizing
dietary quality through retail purchase patterns---a proxy that is
continuously generated, population-wide, and does not require active
participant engagement.

\subsection{OTC Medications During Pregnancy and ASD}

The prenatal medication signal has received increasing attention.
Friel \etal~\cite{friel2021prenatal} found no significant association
between prenatal acetaminophen use and offspring ASD, ADHD, or
intellectual disability in a large population-based analysis,
consistent with a recent comprehensive review~\cite{worringer2024acetaminophen}.
In contrast, Schmidt \etal~\cite{schmidt2019prenatal} reported that
prenatal vitamin use in the first month of pregnancy was associated
with lower ASD recurrence risk and milder symptom severity in the
MARBLES prospective high-risk sibling cohort.

Supplement adherence patterns---particularly folic acid, prenatal
multivitamins, and omega-3 supplements---are partially observable
in retail pharmacy purchase records. Instacart aisle 47
(vitamins/supplements) and aisle 11 (cold/flu/allergy) provide
proxies for supplement adherence and medication risk index,
respectively. While purchase records cannot distinguish between
folic acid and other vitamin formulations at the SKU level without
additional product metadata, aisle-level purchase frequency provides
a meaningful population-scale signal.

\subsection{Socioeconomic Factors and Food Access}

Socioeconomic status and neighborhood food access are established
determinants of prenatal dietary quality. The USDA defines a food
desert as a census tract with both low income (poverty rate $\geq$20\%
or median family income $\leq$80\% of area median) and low food access
(more than 1 mile from a supermarket in urban areas, or more than 10
miles in rural areas)~\cite{usdafara2019}. Panjwani \etal~\cite{panjwani2021food}
demonstrated that food insecurity was associated with worse eating
behaviors and reduced access to preferred foods among families of
children with ASD during COVID-19, disproportionately affecting
low-income households.

Maga\~{n}a \etal~\cite{magana2023racial} found that food insecurity
mediated racial/ethnic disparities in obesity among boys with ASD in
a community sample, with the Area Deprivation Index (ADI) explaining
a significant portion of the disparity. DeVilbiss~\cite{devilbiss2017prenatal}
showed that controlled direct effect analyses in MoBa/ALSPAC indicated
that healthy prenatal dietary patterns explained a modest but
statistically significant portion of socioeconomic disparities in
autism outcomes.

These findings motivate the inclusion of the USDA FARA food desert
score and NSCH SES indicators as a third signal in PrenatalSentinel.
Critically, no prior work has used food desert indices as an explicit
feature in a machine learning model for ASD risk prediction.

\subsection{Federated Learning for Healthcare Prediction}

McMahan \etal~\cite{mcmahan2017communication} introduced the FedAvg
algorithm for federated learning, enabling model training across
distributed clients without sharing raw data. Dwork and
Roth~\cite{dwork2014algorithmic} formalized $(\eps,\delta)$-differential
privacy, providing the theoretical foundation for privacy-preserving
machine learning. The combination of FL and DP has been applied to
a range of healthcare prediction tasks.

Mohammadifar \etal~\cite{mohammadifar2023federated} proposed a
support vector classifier-based federated learning system for ASD
detection across four pooled datasets, reporting accuracy of
approximately 99\% within the study. Alshammari \etal~\cite{alshammari2024explainable}
developed an explainable federated learning framework for toddler
ASD detection using autoencoders, $k$-nearest neighbors, and
multilayer perceptrons with differential privacy integration and
LIME-based explainability. Islam and Khan~\cite{islam2025multimodal}
proposed a federated edge computing architecture for multimodal
autism tracking using speech, motion, and physiological signals,
reporting accuracy and latency improvements on synthetic datasets.

A retail transaction screening preprint~\cite{amebleh2025privacy}
proposed a privacy-preserving consumer analytics framework using
federated gradient aggregation and $k$-anonymity for multistate
telemedicine and consumer behavior data, demonstrating the feasibility
of applying FL to retail purchase records for health-related
prediction.

\textbf{Research gap}: No existing system combines grocery purchase
signals, OTC supplement purchase patterns, and food desert index data
in a unified federated learning framework for prenatal ASD risk
prediction. PrenatalSentinel fills this gap.

% ============================================================
\section{Datasets}
\label{sec:datasets}

\subsection{Signal 1 --- Grocery Food Purchases: Instacart Market Basket}

The Instacart Market Basket Analysis dataset~\cite{instacart2017} is
a publicly available, CC0-licensed dataset containing 3.4 million
grocery orders placed by 206,209 users across 49,688 products
organized into 134 aisles. We use aisle-level purchase frequency
to construct six food-quality scores relevant to prenatal nutrition:
\textit{folate\_score} (leafy greens, legumes, fortified cereals),
\textit{omega3\_score} (packaged seafood, refrigerated seafood),
\textit{iron\_score} (red meat, beans, spinach),
\textit{vitD\_score} (dairy, eggs, fortified products),
\textit{processed\_food\_index} (frozen meals, bakery/desserts,
soft drinks), and \textit{dietary\_diversity\_score} (number of
distinct food aisles purchased per month).

Each score is computed as a normalized purchase frequency ratio
within the relevant aisle cluster, clipped to $[0,1]$. The
processed food index is inverted so that higher values indicate
greater processed food consumption. This operationalization follows
the food-as-medicine framework established in our prior
work~\cite{poreddy2024foodmedicine}.

\subsection{Signal 2 --- OTC Medicine Purchases}

Instacart aisle 47 (vitamins and supplements) and aisle 11
(cold, flu, and allergy) provide proxies for prenatal supplement
adherence and medication risk, respectively. We construct two
OTC features: \textit{supplement\_adherence\_score} (normalized
purchase frequency of aisle 47 products, reflecting prenatal
vitamin, folic acid, and omega-3 supplement purchasing behavior)
and \textit{medication\_risk\_index} (normalized frequency of
aisle 11 products, reflecting over-the-counter analgesic and
antihistamine use). Population-level OTC use rates from the
National Ambulatory Medical Care Survey (NAMCS/NHAMCS)~\cite{cdc2023namcs}
are used to calibrate the synthetic distributions.

\subsection{Signal 3 --- Food Desert Index and SES: USDA FARA + NSCH}

The USDA Food Access Research Atlas (FARA) 2019~\cite{usdafara2019}
provides census-tract-level food access indicators including
\textit{LowAccess}, \textit{LowIncome}, \textit{LowAccessLowIncome},
poverty rate, and median family income. We derive a composite
\textit{food\_desert\_score} from the LAPOP1 and LALOWI1 indicators,
normalized to $[0,1]$ with higher values indicating worse food access.

The National Survey of Children's Health (NSCH) 2020--2022~\cite{nsch2022}
provides ASD diagnosis flags (\texttt{K2Q35A}), maternal health
variables, insurance status, and SES indicators for approximately
50,000 children per survey year. NSCH SES quintiles (based on
family income relative to the federal poverty level) are used to
stratify the \textit{ses\_quintile} feature. Zip-code to census-tract
linkage is performed using the HUD USPS Crosswalk file~\cite{hud2023crosswalk}.

\subsection{Dataset Summary}

Table~\ref{tab:datasets} summarizes the four publicly available
datasets used in this study. All datasets are available under
open licenses or as public-use files, and no individual-level
health data is shared across nodes in the federated architecture.

\begin{table}[t]
\centering
\caption{Datasets Used in PrenatalSentinel}
\label{tab:datasets}
\renewcommand{\arraystretch}{1.2}
\begin{tabular}{p{2.1cm} p{2.0cm} p{1.4cm} p{1.4cm}}
\toprule
\textbf{Dataset} & \textbf{Signal} & \textbf{Size} & \textbf{License} \\
\midrule
Instacart 2017~\cite{instacart2017} & Food + OTC (F, M) & 713 MB & CC0 \\
USDA FARA 2019~\cite{usdafara2019} & Food Desert (E) & 82 MB & Public \\
NSCH 2020--22~\cite{nsch2022} & ASD Labels + SES & $\sim$50K/yr & Public \\
Kaggle ASD~\cite{thabtah2019autism} & ASD Labels & 1,054 rec. & CC0 \\
\bottomrule
\end{tabular}
\end{table}

% ============================================================
\section{Methodology}
\label{sec:method}

\subsection{Feature Engineering Pipeline}

Algorithm~\ref{alg:features} describes the feature extraction
pipeline. For each user $u$ in the prenatal window (defined as
the 9-month period before the estimated delivery date), we compute
a 13-dimensional feature vector $\mathbf{x}_u \in [0,1]^{13}$.
The six food-quality features (F-signal) are derived from
aisle-level purchase frequencies; two OTC features (M-signal)
from supplement and medication aisles; and three SES/food desert
features (E-signal) from USDA FARA and NSCH. Two behavioral
features (basket size and order frequency) complete the vector.

\begin{algorithm}[t]
\caption{PrenatalSentinel Feature Extraction}
\label{alg:features}
\begin{algorithmic}[1]
\REQUIRE Orders $\mathcal{O}_u$ for user $u$ in prenatal window $W$
\REQUIRE USDA FARA tract data $\mathcal{F}$, NSCH SES data $\mathcal{S}$
\ENSURE Feature vector $\mathbf{x}_u \in [0,1]^{13}$
\STATE Compute aisle purchase frequency $a_k = |\{o \in \mathcal{O}_u : \text{aisle}(o)=k\}| / |\mathcal{O}_u|$
\STATE $\text{folate\_score} \leftarrow \text{Normalize}(\sum_{k \in K_{\text{folate}}} a_k)$
\STATE $\text{omega3\_score} \leftarrow \text{Normalize}(\sum_{k \in K_{\omega3}} a_k)$
\STATE $\text{iron\_score} \leftarrow \text{Normalize}(\sum_{k \in K_{\text{iron}}} a_k)$
\STATE $\text{vitD\_score} \leftarrow \text{Normalize}(\sum_{k \in K_{\text{vitD}}} a_k)$
\STATE $\text{proc\_food\_idx} \leftarrow 1 - \text{Normalize}(\sum_{k \in K_{\text{proc}}} a_k)$
\STATE $\text{diet\_div\_score} \leftarrow \text{Normalize}(|\{k : a_k > 0\}|)$
\STATE $\text{supp\_adh} \leftarrow \text{Normalize}(a_{47})$ \COMMENT{Aisle 47: vitamins}
\STATE $\text{med\_risk} \leftarrow \text{Normalize}(a_{11})$ \COMMENT{Aisle 11: cold/flu}
\STATE Link $u$ to census tract via zip code; retrieve $\mathcal{F}[\text{tract}]$
\STATE $\text{food\_desert\_score} \leftarrow \text{Normalize}(\text{LAPOP1} + \text{LALOWI1})$
\STATE $\text{ses\_quintile} \leftarrow \mathcal{S}[\text{income\_fpl\_ratio}]$
\STATE $\text{urban\_flag} \leftarrow \mathbb{1}[\text{urban\_rural\_code} \leq 3]$
\STATE $\mathbf{x}_u \leftarrow [\text{folate}, \omega3, \text{iron}, \text{vitD}, \text{pfi}, \text{dds}, \text{sa}, \text{mr}, \text{fd}, \text{ses}, \text{uf}, \text{of}, \text{bsm}]$
\RETURN $\mathbf{x}_u$
\end{algorithmic}
\end{algorithm}

\subsection{Federated Learning with Differential Privacy}

We adopt a horizontal federated learning architecture with $K=5$
retail nodes (clients), each holding a disjoint partition of the
user base. The global model is a logistic regression classifier
trained via federated gradient averaging (FedAvg)~\cite{mcmahan2017communication}
over $T=40$ communication rounds.

\begin{definition}[$(\eps,\delta)$-Differential Privacy~\cite{dwork2014algorithmic}]
A randomized mechanism $\mathcal{M}: \mathcal{D} \rightarrow \mathcal{R}$
is $(\eps,\delta)$-differentially private if for all adjacent datasets
$D, D' \in \mathcal{D}$ differing in one record, and all $S \subseteq \mathcal{R}$:
\[
\Pr[\mathcal{M}(D) \in S] \leq e^\eps \Pr[\mathcal{M}(D') \in S] + \delta.
\]
\end{definition}

At each client $k$ in round $t$, local gradients are clipped to
$\ell_2$ norm $C=1.0$ and Gaussian noise is added before transmission:
\begin{equation}
\tilde{g}_k^{(t)} = \text{clip}(g_k^{(t)}, C) + \mathcal{N}\!\left(0,\, \frac{C^2 \cdot 2\ln(1.25/\delta)}{\eps^2 n_k^2}\right)
\label{eq:dp_noise}
\end{equation}
where $n_k$ is the number of samples at client $k$ and $\delta=10^{-5}$.
The global model is updated as:
\begin{equation}
w^{(t+1)} = w^{(t)} - \eta_t \cdot \frac{1}{K}\sum_{k=1}^{K} \tilde{g}_k^{(t)}
\label{eq:fedavg}
\end{equation}
with learning rate $\eta_t = \eta_0 / (1 + 0.03t)$ and $\eta_0=0.10$.

\subsection{Equity-Aware Evaluation}

Following the responsible AI framework for health prediction, we
evaluate PrenatalSentinel across five SES quintiles (Q1=lowest to
Q5=highest) and two food desert subgroups (high: food desert score
$>0.5$; low: $\leq 0.5$). The primary equity metric is the
\textit{AUC gap}: the difference in AUC-ROC between the highest-
and lowest-performing subgroups. A smaller gap indicates more
equitable performance.

\subsection{Signal Ablation Protocol}

To quantify the marginal contribution of each signal, we train
the gradient boosting classifier (100 trees, depth 3, learning
rate 0.07) under five configurations: F-only (6 features),
F+M (8 features), F+M+E (11 features), F+M+E+B (13 features),
and All-13. All configurations use 5-fold stratified cross-validation.

% ============================================================
\section{Experimental Setup}
\label{sec:experiments}

\subsection{Synthetic Cohort Construction}

Since direct linkage between individual Instacart purchase records
and NSCH ASD diagnoses is not available in the public domain, we
construct a synthetic case-enriched research cohort of $N=5{,}000$
records. This approach follows established practice in federated
learning feasibility studies~\cite{mohammadifar2023federated,
alshammari2024explainable} and is transparently disclosed per
JBHI data citation requirements.

The synthetic cohort is generated via a logistic outcome model
calibrated to known epidemiological associations:
\begin{equation}
\log\!\left(\frac{p_{\text{ASD}}}{1-p_{\text{ASD}}}\right) = \beta_0 + \boldsymbol{\beta}_F^\top \mathbf{x}_F + \boldsymbol{\beta}_M^\top \mathbf{x}_M + \boldsymbol{\beta}_E^\top \mathbf{x}_E + \epsilon
\label{eq:outcome}
\end{equation}
where $\boldsymbol{\beta}_F$ encodes protective effects of folate
($-1.50$), omega-3 ($-1.00$), iron ($-0.80$), and vitamin D ($-0.70$)
consistent with the literature~\cite{guo2020prenatal, lyall2021maternal};
$\boldsymbol{\beta}_M$ encodes supplement adherence ($-1.20$) and
medication risk ($+0.90$); $\boldsymbol{\beta}_E$ encodes food desert
score ($+0.75$) and SES ($-0.50$); and $\epsilon \sim \mathcal{N}(0,1.0)$
represents unexplained variance. The oracle AUC of the true probability
model is 0.799, and the synthetic cohort contains 386 ASD cases
(7.7\% prevalence), consistent with a high-risk enriched screening
population.

\subsection{Baseline Classifiers}

We evaluate three centralized baseline classifiers using 5-fold
stratified cross-validation: logistic regression (LR, $C=1.0$,
class-balanced), random forest (RF, 150 trees, max depth 10), and
gradient boosting (GB, 150 trees, depth 4, learning rate 0.06).
All models use standardized features.

\subsection{Federated Learning Configuration}

The federated model uses $K=5$ clients, $T=40$ rounds, and is
evaluated at four privacy budgets: $\eps \in \{1, 5, 10, \infty\}$
with $\delta=10^{-5}$. The no-DP case ($\eps=\infty$) serves as
the federated performance ceiling.

\subsection{Evaluation Metrics}

Primary metric: AUC-ROC (area under the receiver operating
characteristic curve). Secondary metrics: F1-score, precision,
recall, and average precision (AP). Equity metrics: AUC by SES
quintile and food desert subgroup.

% ============================================================
\section{Results}
\label{sec:results}

\subsection{Baseline Model Performance}

Table~\ref{tab:baseline} presents the 5-fold cross-validated
performance of the three baseline classifiers. Logistic regression
achieves the highest AUC of 0.672, consistent with the linear
separability of the logistic outcome model. Random forest (AUC
0.642) and gradient boosting (AUC 0.634) show slightly lower AUC
but comparable precision, reflecting the challenge of class
imbalance (7.7\% positive rate). The oracle AUC of 0.799 confirms
that the synthetic cohort contains a meaningful but realistically
noisy signal.

\begin{table}[t]
\centering
\caption{Baseline Classifier Performance (5-fold CV, $N=5{,}000$)}
\label{tab:baseline}
\renewcommand{\arraystretch}{1.2}
\begin{tabular}{lcccc}
\toprule
\textbf{Model} & \textbf{AUC} & \textbf{F1} & \textbf{Precision} & \textbf{Recall} \\
\midrule
Logistic Regression & 0.672 & 0.201 & 0.121 & 0.611 \\
Random Forest       & 0.642 & 0.106 & 0.148 & 0.083 \\
Gradient Boosting   & 0.634 & 0.010 & 0.143 & 0.005 \\
\midrule
\textit{Oracle (true prob.)} & \textit{0.799} & --- & --- & --- \\
\bottomrule
\end{tabular}
\end{table}

\subsection{Federated Learning with Differential Privacy}

Table~\ref{tab:fl} presents the federated learning results across
four privacy budgets. The federated model achieves AUC 0.685 at
$\eps=1$ (strict privacy), compared to the no-DP ceiling of
AUC 0.685. The privacy-utility tradeoff is remarkably flat across
the evaluated range ($\eps \in \{1, 5, 10, \infty\}$), indicating
that the Gaussian noise mechanism at the gradient level does not
significantly degrade model performance for this feature space and
client configuration ($K=5$, $n_k \approx 1{,}000$).

\begin{table}[t]
\centering
\caption{Federated Learning Performance vs.\ Privacy Budget}
\label{tab:fl}
\renewcommand{\arraystretch}{1.2}
\begin{tabular}{lcccc}
\toprule
\textbf{Privacy ($\eps$)} & \textbf{AUC} & \textbf{$\Delta$AUC vs.\ LR} & \textbf{Clients} & \textbf{Rounds} \\
\midrule
$\eps=1$ (Strict)   & 0.685 & $+0.013$ & 5 & 40 \\
$\eps=5$            & 0.685 & $+0.013$ & 5 & 40 \\
$\eps=10$           & 0.685 & $+0.013$ & 5 & 40 \\
$\eps=\infty$ (No DP) & 0.685 & $+0.013$ & 5 & 40 \\
\midrule
\textit{LR Baseline} & \textit{0.672} & --- & --- & --- \\
\bottomrule
\end{tabular}
\end{table}

\begin{figure}[t]
  \centering
  \includegraphics[width=\columnwidth]{fig2_privacy_utility.png}
  \caption{Privacy-utility tradeoff: AUC-ROC vs.\ privacy budget $\eps$.
           PrenatalSentinel maintains AUC~$\approx 0.685$ across the full
           range $\eps \in [0.5, \infty)$, demonstrating negligible
           privacy-utility degradation under strict differential privacy.}
  \label{fig:pu}
\end{figure}

The flat privacy-utility curve (Fig.~\ref{fig:pu}) demonstrates
that PrenatalSentinel can operate at strong privacy guarantees
($\eps=1$, $\delta=10^{-5}$) without sacrificing predictive
performance relative to the no-DP baseline. This finding is
consistent with the theoretical analysis of the Gaussian mechanism
under large client sample sizes ($n_k \sim 10^3$), where the
per-gradient noise magnitude is small relative to the gradient
signal.

\subsection{Equity Audit}

Table~\ref{tab:equity} presents the AUC-ROC by SES quintile and
food desert subgroup. There is a non-monotonic pattern across
quintiles, with Q3 (median SES) achieving the highest AUC (0.588)
and Q5 (highest SES) and Q1 (lowest SES) achieving the lowest
AUC values (0.518 and 0.528, respectively). This U-shaped pattern
reflects the SES-stratified ASD prevalence in the synthetic cohort:
Q1 has the highest ASD rate (13.1\%) but the weakest signal
separability due to confounding by food desert effects, while Q5
has the lowest ASD rate (2.8\%) with insufficient positive cases
for robust discrimination.

\begin{table}[t]
\centering
\caption{Equity Audit: AUC-ROC by SES Quintile and Food Desert Subgroup}
\label{tab:equity}
\renewcommand{\arraystretch}{1.2}
\begin{tabular}{lcccc}
\toprule
\textbf{Subgroup} & \textbf{AUC} & \textbf{$n$} & \textbf{$n_{\text{ASD}}$} & \textbf{Mean FD Score} \\
\midrule
Q1 (Lowest SES) & 0.528 & 1,030 & 135 & 0.307 \\
Q2              & 0.573 & 974   & 95  & 0.251 \\
Q3 (Median)     & 0.588 & 1,028 & 73  & 0.190 \\
Q4              & 0.569 & 974   & 55  & 0.133 \\
Q5 (Highest SES)& 0.518 & 994   & 28  & 0.068 \\
\midrule
High Food Desert ($>0.5$) & 0.607 & 237   & 33  & --- \\
Low Food Desert ($\leq 0.5$)  & 0.628 & 4,763 & 353 & --- \\
\midrule
\textit{AUC Gap (High vs.\ Low FD)} & \multicolumn{4}{c}{\textit{0.021 (3.3\% relative)}} \\
\bottomrule
\end{tabular}
\end{table}

The food desert subgroup analysis reveals that high food desert
populations (food desert score $>0.5$) achieve AUC 0.607 compared
to 0.628 for low food desert populations---a 3.3\% relative gap.
This quantifies the AI detection disadvantage for the most
food-insecure populations, who are simultaneously at higher ASD
risk and harder to screen accurately. Fig.~\ref{fig:equity}
visualizes this disparity alongside the food desert score gradient
across SES quintiles.

\begin{figure}[t]
  \centering
  \includegraphics[width=\columnwidth]{fig3_equity_audit.png}
  \caption{Equity audit: (left) AUC-ROC by SES quintile with mean AUC
           reference line; (right) mean food desert score by SES quintile.
           The AUC gap between high and low food desert populations is
           3.3\% (0.607 vs.\ 0.628), quantifying the AI detection
           disadvantage for the most food-insecure families.}
  \label{fig:equity}
\end{figure}

\subsection{Signal Ablation Study}

Table~\ref{tab:ablation} presents the ablation results. The
F+M+E configuration (food + OTC + SES/food desert, 11 features)
achieves the highest AUC of 0.650, a 2.4\% improvement over
food signals alone (F-only: 0.635). Adding the SES and food
desert features (E-signal) contributes 1.7 AUC points beyond
the F+M configuration (0.633), confirming the independent
predictive value of neighborhood food access. The basket-level
behavioral features (order frequency, basket size) do not provide
additional discriminative power beyond the F+M+E configuration
in this synthetic cohort.

\begin{table}[t]
\centering
\caption{Signal Ablation Study (5-fold CV, Gradient Boosting)}
\label{tab:ablation}
\renewcommand{\arraystretch}{1.2}
\begin{tabular}{lccc}
\toprule
\textbf{Configuration} & \textbf{Features} & \textbf{AUC} & \textbf{$\Delta$AUC vs.\ F-only} \\
\midrule
F only (Food)     & 6  & 0.635 & --- \\
F+M (Food+OTC)    & 8  & 0.633 & $-0.002$ \\
F+M+E (Full)      & 11 & 0.650 & $+0.015$ \\
F+M+E+B (+Basket) & 13 & 0.646 & $+0.011$ \\
All 13 Features   & 13 & 0.646 & $+0.011$ \\
\bottomrule
\end{tabular}
\end{table}

Feature importance analysis (gradient boosting) identifies
\textit{processed\_food\_index} (0.113) and
\textit{food\_desert\_score} (0.107) as the two most important
features, followed by \textit{iron\_score} (0.106),
\textit{omega3\_score} (0.103), and \textit{folate\_score} (0.091).
This ranking is consistent with the epidemiological literature,
which identifies dietary quality and food access as the strongest
modifiable prenatal risk factors.

\begin{figure}[t]
  \centering
  \includegraphics[width=\columnwidth]{fig4_ablation.png}
  \caption{Signal ablation study: AUC-ROC by feature configuration.
           The F+M+E configuration (food + OTC + SES/food desert,
           11 features) achieves the highest AUC of 0.650, a 2.4\%
           improvement over food signals alone (F-only: 0.635).
           Dashed line indicates logistic regression baseline (0.672).}
  \label{fig:ablation}
\end{figure}

\begin{figure}[t]
  \centering
  \includegraphics[width=\columnwidth]{fig5_roc_curves.png}
  \caption{ROC curves for (left) baseline classifiers and (right)
           federated learning models at four privacy budgets.
           FL models consistently outperform the LR baseline.
           All FL curves are nearly indistinguishable, confirming
           the flat privacy-utility tradeoff.}
  \label{fig:roc}
\end{figure}

% ============================================================
\section{Discussion}
\label{sec:discussion}

\subsection{Interpretation of Results}

The federated learning model (AUC 0.685 at $\eps=1$) outperforms
the best centralized baseline (LR, AUC 0.672) by 1.3 AUC points.
This is counterintuitive at first glance---federated models with
DP noise would be expected to underperform centralized models due
to gradient perturbation. The improvement arises from the implicit
regularization effect of federated gradient averaging across five
heterogeneous client partitions, which reduces overfitting to the
training distribution. This finding aligns with theoretical analyses
showing that FedAvg can outperform centralized gradient descent
when client data distributions are heterogeneous~\cite{mcmahan2017communication}.

The flat privacy-utility curve across $\eps \in \{1, 5, 10, \infty\}$
is a particularly noteworthy finding. It suggests that for this
feature space and client configuration, the privacy cost of strict
DP ($\eps=1$) is negligible. This is consistent with the large
per-client sample size ($n_k \approx 1{,}000$), which dilutes the
per-gradient noise magnitude. In practice, retail loyalty programs
with millions of users per node would exhibit even flatter
privacy-utility curves.

The equity audit finding---that high food desert populations achieve
lower AUC than low food desert populations (0.607 vs.\ 0.628)---is
a critical result for public health policy. It implies that a
PrenatalSentinel-style deployment would be least effective precisely
for the populations most in need of early intervention. This
motivates two future directions: (1)~targeted data collection in
food desert zip codes to improve model calibration for low-SES
populations, and (2)~fairness-constrained training that explicitly
optimizes for equity across SES strata.

\subsection{Limitations}

Several limitations of this study must be acknowledged. First, the
synthetic cohort construction, while calibrated to known
epidemiological associations, does not capture the full complexity
of real-world confounding. The outcome model assumes linear additive
effects and may underestimate interaction effects between dietary
quality and genetic risk factors. Second, the Instacart dataset does
not include gestational timestamps, requiring proxy estimation of
the prenatal window via order frequency clustering. Third, the
linkage between retail zip codes and USDA FARA census tracts
introduces measurement error due to the ecological fallacy.
Fourth, purchase records cannot distinguish between individual
household members, and the assumption that prenatal supplement
purchases reflect maternal behavior is an approximation.

\subsection{Pathway to Real-World Deployment}

PrenatalSentinel is designed as a \textit{feasibility framework}
rather than a clinical deployment system. The pathway to real-world
implementation would require: (1)~IRB approval for a prospective
study linking anonymized loyalty card data to birth registry records,
(2)~partnership with one or more major grocery retailers, (3)~validation
on a real longitudinal cohort with confirmed ASD diagnoses,
(4)~regulatory review under HIPAA safe harbor provisions for retail
data, and (5)~clinical integration with pediatricians for risk
communication and referral. The IACC Strategic Plan 2021--2023
explicitly calls for novel data sources for ASD surveillance,
providing a policy mandate for this research direction.

% ============================================================
\section{Responsible AI and Ethics}
\label{sec:ethics}

\subsection{Data Ethics and Privacy}

All datasets used in this study are publicly available under open
licenses or as public-use files, as required by JBHI data citation
policy. The Instacart dataset is released under CC0 (no rights
reserved). The USDA FARA and NSCH datasets are federal public-use
files with no individual identifiers. The Kaggle ASD screening
dataset is released under CC0. No individual-level health records
are shared across federated nodes in the proposed architecture.

The $(\eps,\delta)$-DP guarantee with $\eps=1$, $\delta=10^{-5}$
provides strong formal privacy protection. Under the Gaussian
mechanism, the probability that any individual's data can be
distinguished from a neighboring dataset is bounded by $e^1 \approx 2.72$
with probability $1-10^{-5}$.

\subsection{Algorithmic Fairness}

The equity audit presented in Section~\ref{sec:results} quantifies
the AI detection gap across SES quintiles and food desert subgroups.
We acknowledge that the current model exhibits lower performance
for high food desert populations, and we propose fairness-constrained
training as a future direction. Deployment of PrenatalSentinel
without addressing this equity gap would risk exacerbating existing
health disparities.

\subsection{Synthetic Data Disclosure}

In accordance with responsible AI practices and JBHI transparency
requirements, we explicitly disclose that the experimental results
in this paper are derived from a synthetic cohort calibrated to
published epidemiological associations. The synthetic cohort is
not a substitute for real-world validation. We release the data
generation code and parameters at the project repository to enable
reproducibility and community validation.

\subsection{Intended Use and Misuse Prevention}

PrenatalSentinel is intended as a population-level public health
surveillance tool, not as an individual clinical diagnostic. Risk
scores should not be communicated to individual families without
clinical mediation. The system is not designed to predict ASD with
certainty; it identifies population-level risk gradients to guide
resource allocation for early screening programs.

% ============================================================
\section{Conclusion}
\label{sec:conclusion}

We presented PrenatalSentinel, a privacy-preserving federated
learning framework for prenatal ASD risk prediction from multimodal
retail purchase signals. Using publicly available datasets---Instacart
Market Basket, USDA Food Access Research Atlas, NSCH, and Kaggle
ASD screening---we demonstrated that a federated gradient aggregation
model with $(\eps=1,\delta=10^{-5})$-differential privacy achieves
AUC 0.685, outperforming centralized baselines (best AUC 0.672)
while providing formal privacy guarantees. A signal ablation study
confirmed that the F+M+E configuration (food + OTC + SES/food desert)
achieves the highest AUC of 0.650, validating the multimodal
hypothesis. An equity audit quantified a 3.3\% relative AUC gap
between high and low food desert populations, identifying a
critical fairness challenge for future work.

PrenatalSentinel addresses five open research gaps: the prenatal
purchase signal gap, the food-as-medicine operationalization gap,
the prenatal medication signal gap, the privacy-preserving
population surveillance gap, and the equity/disparity gap. With
ASD prevalence now at 1 in 31 children and an economic burden
exceeding \$461 billion annually, scalable, privacy-preserving
prenatal surveillance systems represent a high-impact research
frontier. Future work will focus on real-world validation with
longitudinal cohort data, fairness-constrained federated training,
and clinical integration pathways.

% ── Acknowledgments ───────────────────────────────────────
\section*{Acknowledgment}
The author acknowledges the open-data contributions of Instacart,
the USDA Economic Research Service, the U.S. Census Bureau (NSCH),
and the Kaggle community. This work was conducted independently
without external funding.

% ── References ────────────────────────────────────────────
\bibliographystyle{IEEEtran}
\bibliography{references}

\end{document}
